\chapter{The Electron Is Named}

\section{Naming makes it measurable}

The corridor carries four faces of one residue. This chapter watches the corridor settle, for the first time, into a
\emph{named} object --- and the name is not decoration laid over a reading already complete. The naming is what makes the
reading a reading. An unnamed corridor is a quantity in transit; a named one is a thing the machine can measure again,
point to, and hold steady. This section is about the act of naming, and about the humble combinatorial fact that forces
it: the pigeonhole principle.

Begin with the machine's situation. It has built an endless supply of formal terms --- the tower climbs without bound,
each rung a new construction --- and it has only finitely many boxes to sort them into. When you must place infinitely
many items into finitely many boxes, some box receives more than one; that is the pigeonhole principle, the plainest
counting fact there is, and it is exact rather than approximate. The machine's endless terms, sorted into its finite
boxes, must therefore collide: two distinct terms are forced to share a box. That forced collision is not a failure of
resolution. It is the machine discovering that two of its endless constructions are, for its purposes,
\emph{the same} --- indistinguishable, poolable, one. Finiteness compels the identification, and the identification is a
name.

This is why naming ``makes it measurable'' rather than merely labeling something already measured. Before the pigeonhole
forces a collision, the machine has only terms in transit --- distinct constructions, none of them settled. The
collision is the first place the machine can say: \emph{these are one thing, and here is the box that holds it}. A box
with a forced occupant is a name: a stable address the machine can return to and find the same thing. And a thing with a
stable address is a thing the machine can measure, because measuring requires re-finding --- you cannot take a reading
twice of something you cannot locate twice. The pigeonhole gives the corridor an address; the address is the name; and
the name is what turns a quantity in transit into an object a reading can be taken of. Naming is not cosmetic. It is the
precondition of repeatability, which the last part made the whole of honesty.

The finiteness doing this work is the same finiteness the whole book insists on, and here it pays a dividend rather than
exacting a cost. Elsewhere finiteness is the machine's limitation --- the floor below which it cannot see. Here it is the
machine's power: because the boxes are finitely many, the endless terms \emph{must} collide, and the collision hands the
machine a name for free. An instrument with infinitely many boxes could keep every term apart forever and would never be
forced to name anything; it would have infinite resolution and no objects. This machine has finite resolution and,
exactly because of it, objects --- named, addressed, measurable. Finiteness is what forces the corridor to become a
thing.

In the two verbs naming is a funge the pigeonhole forces. The machine tanged its terms apart, endlessly, up the tower ---
each a distinct construction. But finitely many boxes cannot keep endlessly many terms apart, so the pigeonhole forces a
funge: two terms it had tanged apart are pooled into one box, declared the same. The name is that forced funge --- the
box that must hold more than one term, read as the single object those terms all name. Where the machine's tanging was
free and endless, the funge here is compelled: it did not choose to pool these terms, finiteness made it. And a compelled
funge is exactly a name that is not fiat --- a point the last section of this chapter will make precise.

\section{The stable reading}

The pigeonhole forces a name; this section says what is named. The corridor, settled into its forced box, enters the
book as a specific, stable reading --- the object the physical gauge will call the electron, and which this gauge reads
as the corridor's first fixed point. The reading is stable in the exact sense that matters to a measuring machine: it
returns the same value every time it is read, and it carries the same simple tag --- a single decided step in one
direction, which the machine writes as $-1$.

Take the stability first, because it is the whole of why this reading, and not another, gets a name. Many quantities run
down the corridor; most are in transit, changing from pass to pass, never the same twice. The named reading is the one
that does \emph{not} change --- the value a pass returns unaltered, the fixed point of the descent. When the machine
sends this reading around its loop, it comes back itself; the loop closes on it exactly. That is what makes it nameable:
a reading that reruns to itself is a reading the machine can stand behind, by the repeatability rule of the last part. The
electron, in this gauge, is simply the corridor's stable reading --- the fixed point the loop returns unchanged, the one
value in the flow that holds still long enough to be named and measured.

The tag it carries, $-1$, is the plainest reading the corridor can present, and this gauge reads it as a single decided
step in one direction. It is not a magnitude and not yet a coupling; it is a count of one, signed --- the smallest
nonzero reading, a single unit of the count face resolved to its direction. The physical gauge reads this same $-1$ as a
charge; this gauge reads it as a unit step, one tick of the tally with its orientation decided. The two readings are the
same corridor at the same port, named in two languages, and this book keeps the computational one: the electron enters
as the stable reading whose count face is one step, decided negative. That is all the machine claims about it here --- a
fixed point, tagged with a single decided step, and stable enough to name.

Why $-1$ and not some other stable value is a question this gauge answers only in part, and it is honest about the
limit. The direction --- that the decided step falls one way rather than the other --- the machine reads off the carried
fact, a decidable bit it can settle. The \emph{unit}, that the count is one and not more, the pigeonhole gives: the
smallest forced collision is between two terms, and two terms sharing one box is a count of one. But why this stable
reading sits where it does in the corridor, and not the magnitude of anything, is not the naming's to say. The naming
fixes the object and its count face; the magnitude faces --- the bound, the cost --- are the far chapters', and even
there they arrive as brackets. The stable reading is named here, fully, as what it is: a fixed point with a unit,
decided count. Its size, in the faces that have one, is not yet asked.

In the two verbs the stable reading is the funge the loop certifies. The machine funges the colliding terms into one box
--- the name --- and then tanges the box's reading apart from the flow by the one characteristic that distinguishes it:
it holds still. The stable reading is the funge that reruns to itself, tanged out of the transit by its fixity. And the
tag $-1$ is a second tange within the name --- the decided direction of the single step, selected off the carried fact.
So the electron is a funge named and a tange tagged: the pooled box that holds still, marked with the one decided step it
carries. The last thing the chapter owes is why this naming is forced by counting rather than declared by fiat.

\section{Named by counting, not fiat}

The corridor has a name and a stable reading. This closing section defends the naming against the one charge that would
undo it: that the machine simply \emph{decided} to call this reading the electron, that the name is a fiat dressed as a
discovery. The defense is the chapter's deepest point, and it is a fact about the shape of the tower. The name is not
fiat because the machine could not have declared it --- the naming is \emph{forced} by counting, and forced in a way the
machine had no freedom to arrange.

Recall why the collision happens. The tower of terms is endless, and it cannot be a straight line running off to
infinity with every term forever distinct, because the machine has only finitely many boxes to hold them. An endless
line through finite boxes is impossible; the line must, at some point, return to a box it has used --- it must
\emph{close into a loop}. The naming is exactly this closure: the point where the endless tower, unable to stay a line,
folds back and reuses a box, forcing two terms to share it. The machine did not choose to fold the tower; finiteness
folded it. The name is the fold, and the fold is compelled by counting alone --- endlessly many terms, finitely many
boxes, no escape. This is Kleene's~\cite{kleene1952} fixed point read combinatorially: a total structure over finite
means must recur, and the recurrence is not posited but forced.

The contrast with fiat is exact, and it is the whole of the section's honesty. A fiat name is one the machine could have
withheld or placed elsewhere --- a free choice, answerable to nothing. This name is answerable to the pigeonhole: the
machine could not have withheld it, because the terms \emph{must} collide, and could not have placed it elsewhere,
because the collision falls where counting puts it, not where the machine prefers. Ask the machine what the naming rests
on and the answer is a count, not a decree: so many terms, so many boxes, collision forced. There is no axiom in it, no
oracle, no free selection --- only the plainest counting fact, run to its conclusion. The name is as compelled as $2 + 2$,
and as free of choice.

This is the chapter that is the logical gauge's summit, and it is worth marking why this gauge climbs it without
planting its own flag here. For the volume that reads the machine as a construction, the naming --- a distinction forced
by counting into a measurable, stable object --- is the peak: it is where the whole apparatus of building-by-telling-apart
delivers its first named thing. This gauge honors that peak and reads it in its own terms --- the pigeonhole, the finite
boxes, the tower folded into a loop --- but its own summits lie elsewhere, at the two chapters where the machine turns
its meter on itself. Here it reads the naming faithfully and hands the climax's full weight to the gauge that owns it.
The electron is named, by counting and not by fiat, and this gauge has read that naming as the forced fold of an endless
tower through finite boxes.

In the two verbs the forced naming is the funge that counting compels. All the way up, the machine tanged --- told each
term apart from the last, freely, endlessly. Finiteness turns that free tanging into a compelled funge: the tower cannot
stay tanged apart forever through finite boxes, so counting forces two terms to pool into one, and that forced pooling is
the name. Fiat would be a funge the machine chose; this is a funge the machine could not avoid. The electron is named by
the one funge counting leaves the machine no choice about --- and with a named, stable object in hand, the next chapter
turns the corridor and watches the reading rotate.
